\scanpage{17}
\mspage{\handwritten{16}}

Le théorème de Thom est alors qu'un $x\handwritten{\in}H^{\bullet}(X)$ est dans $S(X)$ sss il est \deleted{\unsure{stable|stabil}} \inserted{fixe} par $G$, i.e. satisfait $A^{\handwritten{+}}\mathbin{.}x=0$. (Pourquoi cette condition est-elle déjà nécessaire ?) U\inserted{\handwritten{n}}e autre façon de le voir est de dire que le foncteur $H^{\bullet}$ considéré comme allant de la catégorie $\underline S$ précédente dans la catégorie des \deleted{vectoriels} représentations de $G$ dans des vectoriels sur $\underline F_2$, est pleinement fidèle. \deleted{Comme $G$ est} Quand on prend l'enveloppe pseudo-abélienne, on trouve une sous-catégorie pleine de la catégorie des représentations de $G$, mais il est douteux qu'on trouve une catégorie abélienne, voire la catégorie de toutes les représentations de $G$ de \unsure{dimensio|dimension} finie sur $\underline F_2$. (Quelle est la structure de $G$, qui est prolisse; d'après Quillen, il serait pro-\unsure{unipotent|umipotent} ....). Ici $G$ bien sûr joue le rôle d'un groupe de Galois motivique, mais contrairement à la situation familière, il semble qu'il soit fortement non réductif (et d'ailleurs le corps de \unsure{défimtion|défintion|définition} dans tout ceci est $\underline F_2$, et non $\underline Q$ ! ).
\margin{\handwritten{\unsure{Oui}, car son \unsure{anneau} affine \unsure{a|est} une \unsure{infinité|suite|système|algèbre} de \unsure{générateurs|polynômes} ! une \unsure{quantité|infinité} d'\deleted{\unsure{indéte}}\unsure{indéterminées}, \deleted{\unsure{limite}} \unsure{sous} \unsure{relatio.|inductive|réunion} d'\unsure{algèbres} \unsure{polynômiales} à \unsure{augmentation|engendrement} \unsure{et} stables par $\Delta$}}% slanted

Quillen\deleted{t} donne une généralisation de la construction d'une catégorie universelle $\underline B$ au cas où on travaille avec des variétés pas nécessairement propres (mais $f_{\bullet}$ n'étant supposé défini que pour $f$ propre), et quand on introduit une condition d'orientation \inserted{stable} sur les \uline{morphismes} envisagés , via un morphisme $H\unsure{\text{---}|\longrightarrow}\mathrm{BO}(\infty)$, où $H$ est un H-espace (commutatif, associatif et unitaire ?): une orientation de $f:X\unsure{\text{---}|\longrightarrow}Y$ sera un relèvement homotopique de $X\longrightarrow\mathrm{BO}$, défini par \deleted{$\unsure{f^{\bullet}(T_X)|f^{*}(T_X)}$} $T_X-f^{\bullet}(T_Y)$, en $X\longrightarrow H$. les éléments de $H(X,Y)$ seront définis par des diagrammes de morphismes orientés $Z\longrightarrow X,\ Z\unsure{\text{---}|\longrightarrow}Y$, le deuxième étant propre, modulo bords de tels. En tant qu'ensemble, il se définit aussi en termes de la vari\handwritten{é}té $X\mathbin{x}Y$ munie simplement de la famille des supports formée des ensembles fermés propres sur $Y$ .....

\handwritten{Les $B^H_{\unsure{m|n}}(X)$ ($X$ compact disons) peuvent se définir homotopiquement à la Thom, comme les groupes d'homotopie stables d'\unsure{espaces|appl.} $\underline{\mathrm{Hom}}(X,\mathrm M\mathrm H(N-\unsure{m|n}))$ -----.}

\margin{\handwritten{cf exposé Karoubi à \uline{Bourbaki}}}

% archive sheet number: 16
